\documentclass[12pt]{elsarticle}
\usepackage[brazil]{babel}
\usepackage[utf8]{inputenc}
\usepackage{amsmath,amssymb}
\usepackage{enumitem}
\usepackage{geometry}
\geometry{a4paper, margin=2.5cm}
\usepackage{caption} % para usar \caption*
\makeatletter
\def\ps@pprintTitle{%
	\let\@oddhead\@empty
	\let\@evenhead\@empty
	\def\@oddfoot{\footnotesize \textit{Submitted by Jonatha} Costa\hfill}%
	\let\@evenfoot\@oddfoot
}

\usepackage{url, geometry, hyperref,enumitem}
\usepackage[abnt-emphasize=bf,abnt-and-type=e,alf]{abntex2cite} % Citações ABNT
\usepackage{multicol}   % Multiplas colunas
\geometry{
	left=2.5cm,
	right=2.5cm,
	top=2.5cm,
	bottom=2.5cm
}

\hypersetup{%
	colorlinks = true,
	linkcolor = blue, % Define a cor do texto
	citecolor = blue!85, % Define a cor da citação
}
\begin{document}

	
\begin{frontmatter}
	\title{\bf MÉTODOS NUMÉRICOS\\
	\textnormal{\normalsize Engenharia de Controle e Automação \& Engenharia Mecânica}
}
\author{
\vspace{-0.5cm}
\textnormal{
 	Lista Geral de exercícios\\
	Prof. Jonatha Costa\\
	jonatha.costa@ifce.edu.br}
}
	\begin{abstract} 
Este documento reúne uma lista de exercícios de Métodos Numéricos para os cursos de Engenharia de Controle e Automação e Engenharia Mecânica do IFCE \textit{campus} Maracanaú,
alinhada ao Plano de Unidade Didática (PUD). São trabalhados temas centrais como análise de erros, métodos para equações não lineares, sistemas lineares, interpolação,
ajuste de curvas, integração e diferenciação numérica. Os exercícios são direcionados para aplicações reais de engenharia, com ênfase em modelagem,
simulação e avaliação do desempenho dos algoritmos. O foco é direto: consolidar rigor analítico e formar o aluno para implementar métodos numéricos de forma eficiente, 
visando aplicações computacionais mais avançadas.	
\end{abstract}
	
\end{frontmatter}
	

	
\begin{enumerate}
\section*{Conversão - norma IEEE 754/2008}

	% *****************************************************************
\item Represente os números a seguir nos formatos solicitados abaixo: 
{\bf 81, 66.25, -0.625, 0.533203125}.
    \begin{enumerate}
        \item Binário convencional;
        \item Binário em ponto flutuante;
        \item Cadeia de 32 bits em precisão simples conforme a norma IEEE--754.
        \item Cadeia de 64 bits em precisão dupla conforme a norma IEEE--754.
    \end{enumerate}
    
\item Faça um {\it script} que receba um número decimal {\bf positivo} e converta-o para:
    \begin{enumerate}
        \item Binário convencional;
        \item Binário em ponto flutuante;
        \item Cadeia de 32 bits em precisão simples conforme a norma IEEE--754;
        \item Cadeia de 64 bits em precisão dupla conforme a norma IEEE--754;
    \end{enumerate}
    
\item Faça um {\it script} que receba um número decimal {\bf negativo} e converta-o para:
	\begin{enumerate}
        \item Binário convencional;
        \item Binário em ponto flutuante;
        \item Cadeia de 32 bits em precisão simples conforme a norma IEEE--754;
        \item Cadeia de 64 bits em precisão dupla conforme a norma IEEE--754;
    \end{enumerate}
    
\item Faça um {\it script} que receba um número decimal com {\bf mantissa} (parte não inteira) e converta-o para:
	\begin{enumerate}
        \item Binário convencional;
        \item Binário em ponto flutuante;
        \item Cadeia de 32 bits em precisão simples conforme a norma IEEE--754;
        \item Cadeia de 64 bits em precisão dupla conforme a norma IEEE--754;
    \end{enumerate}

\item Utilizando POO, ajuste o {\it script} da questão anterior para receber o número decimal em quaisquer modos supracitados e convertê-lo para as saídas 
anteriormente exploradas.

\section*{Regressão linear}


\item Com base no seguinte conjunto de dados \[x = \{2,5,6,8,9,13,15\};\quad y = \{7,8,10,11,12,14,15\}\]:
\begin{enumerate}
    \item Utilize a regressão linear por mínimos quadrados para determinar os coeficientes $a$ e $b$ da função:
    \( y = ax + b\) que melhor se ajusta aos dados.

    \item Determine o erro global do ajuste.

    \item Estime o valor de $y$ para $x = 7.25$ utilizando a função ajustada ($p(x)$).

    \item Trace o gráfico contendo:
    \begin{itemize}
        \item os pontos experimentais (dados fornecidos);
        \item os pontos interpolados $x_int$ e $y_int$;
        \item a reta obtida por regressão linear e polinômio.
    \end{itemize}
\end{enumerate}

	\item Atualize o {\it script} da questão anterior utilizando classes (POO) e analize as respostas 
	para o seguintes dados: \(x = \{-7, -5, -1, 0, 2, 5, 6\}, y = \{15, 12, 5, 2, 0, -5, -9\}\).

\item Considere os vetores abaixo:	
	\begin{align*}
		\mathbf{x}  &= [0,\ 10,\ 20,\ 30,\ 40,\ 50,\ 60,\ 70,\ 80,\ 90,\ 100] \\
		\mathbf{y} &= [1.00,\ 1.40,\ 2.00,\ 2.65,\ 3.36,\ 4.00,\ 4.56,\ 5.03,\ 5.40,\ 5.76,\ 6.20]
	\end{align*}
	
	\begin{enumerate}
		\item Utilizando regressão linear pelo métodos dos mínimos quadrados, escreva um {\it script} que 
		 determine os coeficientes  do polinômio $p(x) = ax  + b$.
		\item Com o polinômio $p(x)$ obtido, calcule os valores correspondentes de $y_{int}$ para os seguintes pontos:
		\(x_{\text{int}} = \{25,\ 35,\ 42\}\)		
		\item Com o polinômio $p(x)$, encontre o novo contra-domínio ($NovoY$) e mostre o erro relativo ao modelo($p(x)$);
		\item Trace um gráfico contendo: $(x \times y)$, $(x \times NovoY)$,($x_{int} \times y_{int}$) do item anterior.
	\end{enumerate}
% *****************************************************************	
	\item Atualize o {\it script} da questão anterior utilizando classes (POO) e faça o código gerar valores 
	aleatórios de $\rho \mid \rho  \subset \mathbf{x}$, os quais devem ser exidos no gráfico a cada nova simulação.


\section*{Polinomios Interpolares de Lagrange e Newton}	
% *****************************************************************
	\item Compare os resultados obtidos pelos métodos Lagrange e Newton a partir de:
	\begin{enumerate}
		\item Utilizando os pontos $(x_{\text{int}}, y_{\text{int}})$ encontrados anteriormente, gere o polinômio interpolador de Lagrange, denotado por $p_1(x)$.
		\item Com os mesmos pontos, gere o polinômio interpolador de Newton, denotado por $p_2(x)$.
		\item Utilizando os polinômios $p_1(x)$ e $p_2(x)$, interpole os valores de $y_{inter}$ para:
		\(x_{\text{int2}} = \{20,\ 30,\ 40\}\)
		\item Compare o erro relativo de cada polinomio com os valores já conhecido no domínio $x$ inicial	
	\end{enumerate}

% *****************************************************************
	\item Atualize o {\it script} da questão anterior utilizando classes (POO).

% *****************************************************************
	

% *****************************************************************
\end{enumerate}
	
	
\end{document}
	% \item \textbf{Interpolação Polinomial}
	% \begin{enumerate}
	% 	\item Utilizando os pontos $(x_{\text{int}}, y_{\text{int}})$ encontrados anteriormente, gere o polinômio interpolador de Lagrange, denotado por $p_1(x)$.
	% 	\item Com os mesmos pontos, gere o polinômio interpolador de Newton, denotado por $p_2(x)$.
	% 	\item Utilizando os polinômios $p_1(x)$ e $p_2(x)$, interpole os valores de $y$ para:
	% 	\[
	% 	x_{\text{int2}} = \{20,\ 30,\ 40\}
	% 	\]
	% 	Compare os resultados obtidos por ambos os métodos.
	% \end{enumerate}
	
	% \item \textbf{Integração Numérica}
	% \begin{enumerate}
	% 	\item A partir de um dos polinômios interpoladores obtidos (Lagrange ou Newton), gere um novo conjunto $Y$ correspondente a uma malha linear de valores $X$ no intervalo $[20,\ 40]$.
	% 	\item Com esse novo conjunto $(X, Y)$, calcule a área sob a curva entre os pontos $x = 25$ e $x = 37$ utilizando os seguintes métodos:
	% 	\begin{itemize}
	% 		\item Regra dos Retângulos (especificar número de subintervalos);
	% 		\item Regra dos Trapézios;
	% 		\item Regra de Simpson (simples ou composta, a critério do aluno).
	% 	\end{itemize}
	% 	\item Apresente e compare os resultados obtidos.
	% \end{enumerate}
	
	% \item \textbf{Ajuste por Regressão Polinomial de Grau Arbitrário}
	% \begin{enumerate}
	% 	\item Implemente um algoritmo que teste regressões polinomiais de diferentes graus (1 a 6) para os dados fornecidos.
	% 	\item Para cada grau, calcule o erro quadrático médio (EQM) entre os valores ajustados e os valores reais.
	% 	\item Apresente graficamente os ajustes obtidos e indique qual grau oferece o melhor equilíbrio entre erro e complexidade.
	% \end{enumerate}
	
	% \item \textbf{Derivada Numérica}
	% \begin{enumerate}
	% 	\item Com base nos dados de $x$ e $y$, estime a derivada de $y$ em cada ponto usando:
	% 	\begin{itemize}
	% 		\item Diferença Progressiva (início);
	% 		\item Diferença Regressiva (fim);
	% 		\item Diferença Central (demais pontos).
	% 	\end{itemize}
	% 	\item Construa um gráfico da função original $y(x)$ e do seu perfil de derivada $y'(x)$.
	% 	\item Analise e interprete os intervalos onde a taxa de variação de $y$ aumenta ou diminui.
	% \end{enumerate}
	% \item \textbf{Consolidação do Código}
	% \begin{enumerate}
	% 	\item Crie um único script (ou notebook) organizado contendo todas as funcionalidades implementadas nas questões anteriores:
	% 	\begin{itemize}
	% 		\item Regressão linear, polinomial e interpolação;
	% 		\item Cálculo de integrais numéricas;
	% 		\item Derivadas numéricas;
	% 		\item Análise gráfica;
	% 		\item Relatórios de erro.
	% 	\end{itemize}
	% 	\item O código deve ser modularizado com funções separadas para cada operação e conter comentários claros explicando a lógica de cada etapa.
	% 	\item Opcionalmente, o código pode ser apresentado como um notebook com células explicativas, gráficos e outputs organizados.
	% \end{enumerate}
